\cotefile{157-4}
\scanpage{148}
\mspage{22}
% archive sheet number: 147

de cardinal majoré par
\[
\underbrace{r^{\pi_0}}_{\substack{\text{nb de \unsure{classes|ces}}\\\text{syst. inductifs}\\\text{dans }\underline M^{\pi},\text{ à}\\\text{isom près}}}
\mathbin{\cdot}
\underbrace{s^{r^{\pi_0}}}_{\substack{\text{nb des}\\\text{flèches d'un}\\\text{tel syst. ind.}\\\text{dans }X}}
= s^{r^{\pi_0}}.
\]

Mais il faut maintenant faire varier $I$,
et nous \unsure{porter|plonger} dans la cat des ens. ordonnés
(filtrants etc) de cardinal $\inserted{\lambda=}\deleted{\unsure{I}}r^{\pi_0}\deleted{\unsure{\longrightarrow}}$.
On voit de suite que cet ensemble de classes d'isomorphie (qui peuvent s'interpréter comme un ens.
quotient de \deleted{\unsure{l'ens}} \inserted{un certain} ens. de parties d'un ens.
fixe $I_0$ de card $\lambda$) est de cardinal $\leq 2^\lambda$.
On trouve donc finalement
\[
\operatorname{card}\widetilde{M^r_{/X}}\leq s^{(r^{\pi_0})}\mathbin{\cdot}2^{(r^{\pi_0})}=s^{(r^{\pi_0})},
\]
sous condition que
\[
r\geq\operatorname{Sup}(\pi_1,\unsure{\mathfrak C_0|C_0}),\qquad s\geq\operatorname{Sup}(\unsure{\pi_1|\mathfrak C_1},\mathfrak C_1).
\]

Il faut prouver l'inégalité similaire pour
$\operatorname{card}\widetilde{\underline{\mathrm{Fl}}^{\,r}M}$ (3.9.). Or, \deleted{\unsure{étant entendu}} \inserted{comme}
pour \deleted{\unsure{x}} $A,B$ dans $M^r$, on a si $r\geq\operatorname{Sup}(\pi_1,C_1)$
\[
\operatorname{card}\operatorname{Hom}_{M}(A,B)\leq r^{(r^{\pi_0})}
\]
en vertu de (3.8) où on fait $s=r$, on trouve que
\[
\begin{aligned}
\operatorname{card}\widetilde{\underline{\mathrm{Fl}}^{\unsure{r}}(M^r_{/X})}
\deleted{\underbrace{\vphantom{\mathrm{Fl}}}_{\unsure{\mathrm{Fl}(M^r)}}}
&\leq\left(\operatorname{card}\widetilde{\deleted{\unsure{M}}M^r_{/X}}\right)^2\mathbin{\cdot}r^{(r^{\pi_0})}\\
&\leq s^{(r^{\pi_0})}\mathbin{\cdot}r^{(r^{\pi_0})}=s^{(r^{\pi_0})}\\[-2pt]
&\hspace{6em}\text{si }s\geq r.
\end{aligned}
\]
% slanted
\margin{cf \unsure{corol.|contr.} \unsure{pge|pg} 23}
(On va regarder plus bas le cas $s\leq r\unsure{{}^{*}|{}^{+}}$\unsure{.}). C'est bien
la formule (3.9.)
